
\subsubsection{2.1 Calibration in Neural Networks and Language Models}

Guo et al. (2017) demonstrated that modern neural networks trained with batch normalization and regularization exhibit systematic overconfidence, and proposed temperature scaling as a post-hoc correction. Their Expected Calibration Error (ECE) framework has become the standard evaluation metric for calibration research. Wang (2023) surveyed calibration methods in deep learning and noted that fine-tuning of pre-trained models introduces calibration degradation, an observation that motivates studying the alignment-calibration relationship. Khanmohammadi et al. (2025) confirmed that ECE varies across model families and benchmarks, supporting the use of MMLU-based calibration evaluation.

None of this prior work distinguishes between mechanistically different sources of miscalibration. A correction appropriate for scale inflation (H1) may be ineffective for boundary restructuring (H2).

\subsubsection{2.2 RLHF Alignment and Calibration Degradation}

Xie et al. (2024) demonstrated that RLHF alignment causes significant calibration degradation in LLaMA-2-Chat and proposed Adaptive Temperature Scaling (ATS), which learns an input-dependent temperature from hidden states and achieves 58-82\% ECE reduction. ATS implicitly assumes an H1-type pattern correctable by rescaling. The present study tests this assumption directly.

Li et al. (2024) used the Pythia alignment ladder (SFT/DPO/PPO variants) to study RLHF effects on trustworthiness, finding that more alignment does not guarantee more trust, with approximately 25\% truthfulness degradation relative to SFT on benchmark suites. However, Li et al. did not measure ECE or Brier calibration. The present study extends their causal design to the calibration dimension.

Coste et al. (2023) studied reward overoptimization, establishing a theoretical mechanism by which reward optimization inflates model confidence. Their analysis predicts confidence inflation on existing choices (consistent with H1), but the present data show near-complete answer redistribution (consistent with H2).

The RLHF framework was established by Ouyang et al. (2022), and DPO by Rafailov et al. (2023). The Pythia model family \citep{biderman2023pythia} provides the only publicly available full alignment ladder sharing identical pretraining data and architecture.

\subsubsection{2.3 Verbal Confidence and Framing Effects}

Chhikara et al. (2025) studied calibration using verbally elicited confidence (0-100 scale) and found that distractors can reduce ECE by up to 90\% in some models. Their findings use a different measurement modality (verbal confidence) than the present study (log-probability softmax ECE). The H3 diagnostic employed here tests whether softmax-based calibration degradation exhibits framing sensitivity within the Pythia family.
